\section{学位论文研究方案}

\subsection{研究目标}

第一，通过研读马克思主义经典著作、党和国家重要政策文献以及国内外相关学术成果，明确当代大学生弱连接交往的范围与特征，结合社会主义核心价值观的育人要求与公民道德生活领域的划分，构建考察大学生价值观变化的四维分析框架与评判标准。

第二，通过收集与分析多平台大学生网络文本资料，梳理弱连接交往背景下大学生在个人生活、人际交往、学习劳动与择业、社会公共四个方面的价值观负面偏移表现，剖析浅层交往规则一步步影响青年日常价值判断的微观作用过程与现实成因，揭示其对马克思主义价值观及社会主义核心价值观培育造成的现实冲击。

第三，立足高校思想政治教育主责，结合弱连接交往的特点与影响机制，从校园实体共同体建设、思政课教学引导、家校协同和校社联动等方面提出切实可行的应对策略，为高校在流动社交环境中提升价值观教育实效、落实立德树人根本任务提供参考。

\subsection{研究内容}

本研究紧扣新时代高校立德树人根本任务，以弱连接交往对大学生价值观的影响为核心对象，围绕概念界定与框架构建、现实表现梳理、作用过程与成因剖析、应对策略提出四个环环相扣的层面依次展开。

首先，进行概念界定与分析框架构建。结合社会网络理论与马克思主义关于人的本质是一切社会关系的总和、现实交往方式反作用于思想观念的基本原理\footcite{sizheng2018}，明确本研究中弱连接交往的判定边界：以交往对象可辨认、存在双向实质互动、具有一定延续性为标准，涵盖大学生日常参与的兴趣社群、跨校学习求职互助、校园搭子往来以及平时联系较少的亲属交往，将开放网络空间中陌生人单向浏览、匿名争吵和舆论站队排除在外。在此基础上，承接2019年《新时代公民道德建设实施纲要》关于社会公德、职业道德、家庭美德和个人品德的领域划分\footcite{gangyao2019}，参照马克思主义理论研究和建设工程重点教材《思想道德与法治》《伦理学》的规范体系\footcite{sixiangdaode2023,lunlixue2020}，贯彻习近平总书记“核心价值观，其实就是一种德”的重要论断\footcite{xi2017wenhua}，结合在校大学生脱离原生家庭日常照料、处于社会生产准备期和校园集体生活阶段的实际，将大学生价值观具体划分为四个维度：一是个人生活价值观，涉及物质消费、精神追求、时间安排与自我修养；二是人际交往价值观，涉及与家人亲属、室友、同学、搭子及婚恋对象等具体他人的相处边界、信任承诺与亲情归属；三是学习劳动与择业价值观，涉及专业学习态度、科研投入、劳动意义认同与职业选择标准；四是社会公共价值观，涉及公共规则遵守、志愿公益、理性表达与社会责任担当。对照马克思主义关于人的全面发展、诚实劳动与集体互助的价值立场，以及社会主义核心价值观的培育要求，确立判断上述四类价值观是否发生负面偏移的具体标准。

其次，梳理弱连接交往背景下大学生价值观负面偏移的具体表现。依托多平台网络文本材料，系统考察当代大学生在典型弱连接场景中的日常表达与价值倾向，揭示其与马克思主义价值观及社会主义核心价值观育人要求之间的现实差距。在个人生活方面，部分学生容易产生脱离自身条件的攀比与跟风消费心理，将外在消费水平视作衡量自我价值的尺度，课余时间也更多流向打卡式消遣，持续性的深度阅读与静心修养明显减少。在学习劳动与择业方面，部分学生将专业学习简化为换取绩点和文凭的短期手段，对需要长期积累的基础研学投入不足，在职业规划上容易在追逐短期高回报与消极逃避竞争之间摇摆，脚踏实地的劳动观与奉献精神有所弱化。在人际交往方面，以校园“搭子”交往和亲戚走动减少为典型表现：与同学和室友相处时，部分学生习惯把往来限定在单一功能协作之内，过度强调边界防御而回避长期信任与互相迁就，遇到矛盾倾向于冷处理或直接断联，集体归属感持续流失；在家庭亲属往来中，部分学生以实用尺度衡量血缘亲疏，对日常缺乏直接利益关联的远亲主动减少联系乃至刻意疏离，亲情联结随之淡化。在社会公共方面，部分学生在校园互助群、跨校兴趣群等松散群体中习惯用情绪化的非黑即白表态替代讲理协商，在现实生活中遵守公共规则的主动性和参与志愿服务的热情有所下降，看待社会问题时容易以偏概全。

再次，剖析弱连接交往影响大学生价值观的微观作用过程与深层成因。在微观作用过程上，重点拆解四条相对独立的传导路径，并结合网络文本材料加以印证：第一，片面展示引发的向上比较过程。弱连接对象之间缺乏全面了解，在消费风气与升学就业竞争压力下，彼此呈现的多是光鲜的消费休闲片段，以及保研、上岸、高薪录用等最终成绩，隐藏了背后的现实条件与长期付出。大学生在高频接触中反复陷入向上比较，容易滋生心理落差与焦虑，进而动摇踏实学习的定力，转向短期功利算计或借由物质消费寻求心理补偿。第二，低义务相处规则向日常关系的延伸过程。受市场等价交换观念向日常生活渗透的影响，以搭子社交为代表的浅层交往往往按需匹配、即时结算、随时散场且无需承担后续责任。反复处于这类互动中，部分学生容易把交往对象当成满足单一功能的实用工具，养成凡事算计得失的防御心态，并将这套相处方式带入宿舍生活、社团合作、婚恋经营乃至亲属往来，削弱同伴间的道义信任，加重对亲戚人情往来的厌烦与回避。第三，浅层满足对深度投入的挤占过程。弱连接交往门槛低、反馈快、摩擦少，能够迅速满足临时陪伴和轻松消遣的需求。长期习惯这种省力即时的满足，容易降低学生在深度阅读、自我修养、维系深厚友谊和耐心沟通亲情上持续花心思的意愿，并在异地求学的空间阻隔下进一步加剧亲情疏离。第四，松散群体环境中的情绪放大与责任淡化过程。校园互助群、兴趣群及阶段性公益群等松散群体聚散自由，成员间缺少稳定的线下共同劳动与长远利益联系。个体围绕公共议题交流时缺乏面对面耐心协商的条件，更倾向于依靠简单贴标签和情绪共鸣换取快速附和，久而久之形成偏激看法；同时，低约束环境也容易淡化集体归属感，使部分学生习惯享受公共便利而回避相应义务，削弱规则意识与社会责任担当。在深层成因上，进一步揭示上述现象背后的现实根源：社会分工细化、人口高频流动、市场功利观念渗透以及升学就业竞争加剧，共同催生出大学生规避人际耗损、寻求低成本互助的防御心理，而移动互联网大幅降低了跨界联系的门槛；家庭规模小型化与异地求学造成的日常陪伴空缺，使校园浅层社交成为便捷替代；高校现有思想政治工作仍主要依靠班级、宿舍等固定组织运转，难以及时感知流动社交中的学生思想动态，教育话语也与青年的日常语境存在隔阂。同时，结合不同交往场景与学生现实强关系支持程度的差异，说明弱连接交往负面影响发挥作用的触发条件与适用限度。

最后，提出化解负面影响的高校思想政治教育应对策略。立足高校育人主责，坚持尊重青年合理个人边界与强化主流价值引领相结合，构建四方面协同路径：第一，做实校园实体集体与家校亲情联结。依托党团支部、跨学科科研团队、学生社团和实践项目，创设需要长期合作、共同做事的真实集体场景，促使浅层搭子在并肩解决实际问题的过程中转化为彼此信任的伙伴；同时借助家校协同沟通与返乡实践，引导学生超越功利算计、珍视家庭亲情。第二，反向借用弱连接跨界优势实施分类纠偏。主动搭建在校生与国家重点产业一线建设者、科研骨干及基层工作者的常态化交流渠道，用完整真实的劳动奋斗经历纠偏功利浮躁的学业与择业观；通过校园读书共同体与理性消费研讨引导青年合理规划课余生活；通过规范有序的校园公共事务协商与社区志愿服务，将线上关切转化为线下理性参与和公共担当。第三，强化思政课的方法论训练与日常叙事转化。在思想政治理论课教学中，结合青年熟悉的日常经历讲透马克思主义关于社会关系、诚实劳动与理性消费的道理，指导学生全面看待交往对象的真实处境，自觉辨析把人际关系当工具、盲目攀比等倾向，推动社会主义核心价值观融入青年易于理解与参与的日常实践。第四，拓展校社协同的外部育人空间。联动高校与地方共建的社区实践基地、青年服务站及校友行业导师网络，将稳定的线下协作共同体延伸至校园之外，涵养互信友善的青年交往风尚。

% 【此处保留原有的“毕业论文框架目录”部分】

\subsection{拟采用的研究方法}

本研究坚持辩证唯物主义和历史唯物主义的根本方法，将弱连接交往对大学生价值观的影响置于经济社会变迁与校园现实竞争压力的物质条件中加以考察，既客观承认其在拓宽信息渠道方面的积极作用，也重点剖析其浅层互动规则诱发观念偏差的消极影响。具体采用以下三种研究方法。

\subsubsection{文献研究与理论分析法}

系统精读《马克思恩格斯文集》中的《1844年经济学哲学手稿》《关于费尔巴哈的提纲》《德意志意识形态》《共产党宣言》《资本论》（第一卷）等经典文本，梳理马克思主义关于人的社会关系本质、交往方式反作用于思想观念、劳动观、消费观及人的全面发展的核心论述；研读党的创新理论文献、高校思想政治工作政策文件及马工程重点教材，明确考察大学生四维价值观变化的价值准绳。同时，广泛检索国内外关于弱关系理论、轻社交、搭子交往、青年亲属疏离及大学生价值观演变的代表性研究，立足马克思主义立场对社会学、传播学与思想政治教育学成果加以整合，搭建本研究的理论分析框架。

\subsubsection{网络文本挖掘法}

用于从自然交往场景中收集大学生关于弱连接互动的真实表达，并从宏观层面把握高频话题与话语关联。为避免开放网络社区中用户身份混杂、泛泛而谈的问题，数据采集采取平台分层、高校身份核验与弱连接情境限定相结合的筛选策略：以代表性高校公开论坛作为身份校准样本，重点覆盖大学生活跃的小红书、豆瓣高校及青年小组、知乎高等教育及应届生板块近三年的公开主题帖与互动评论；通过高校专属标签、校内板块属性及发帖者自述的本科生、研究生、室友、保研、春招等在校信息过滤非大学生群体，严格按照“对象可辨认、有实质回应、有一定延续性”的标准，定向收集校园拼单与搭子互动复盘、轻社交边界讨论、减少亲戚往来反思、考研与春招打卡互助（围绕“绩点”“内卷”“上岸”等竞争议题展开）以及校园互助群与公益社群的持续互动帖文，剔除陌生人单向浏览和匿名争吵内容。在遵守平台公开访问规则与匿名脱敏要求的前提下，先用Python脚本辅以人工检索获取公开正文、回复串及点赞评论数据，再经过去重、清除广告、剔除不足10字的极短无意义评论与人工核查灌水噪音，按各板块综合互动量排名前20\%且有效回复不少于10条的标准截断抽样，保留每帖前20条高赞回复，按平台与四类价值观均衡配置，建成涵盖3000至5000条有效主题帖及对应回复的大学生弱连接交往文本数据库。在此基础上，借助Python开展分词、词频统计、关键词重要性测算与词语共现分析，并严格区分用于检索的场景词（如搭子、打卡群、绩点、上岸）与文本中自然浮现的价值态度词（如消费心态、功利焦虑、边界防御、亲属疏离与公共态度），防止因检索词反复出现而产生同义循环；同时运用短文本主题聚类方法归类不同社区的潜在讨论主题，从宏观上呈现弱连接交往场景与四类价值观表达之间的关联结构。

\subsubsection{质性内容分析法}

鉴于大学生在网络社区的发言常常带有反讽、调侃和圈层隐语，单靠计算机自动统计容易误读真实态度，研究在宏观主题归类之后，按平台与主题分布分层抽取400组典型主题帖及其高赞回复作为深度分析样本，导入质性分析软件NVivo，依据四维价值观界定标准编制操作化编码手册，开展人工精细编码与话语分析。为保证判断标准一致可靠，先抽取15\%的样本由两名研究者独立试编码，检验编码一致性并就分歧条目讨论修订细则，确认达标后再完成全部样本编码，同时辅以跨平台比对以减少单一社区偏差。通过将人工深度编码还原的微观比较心理、交往规则延伸、浅层满足替代与公共责任淡化线索，同宏观统计结果相互印证，为微观作用过程分析提供扎实支撑。研究严守网络文本资料的解释限度，立足话语线索辨析与作用过程阐释，不作超出数据支撑的统计因果推论。

\subsection{研究的可行性分析}

第一，理论与文献基础扎实。马克思主义经典著作关于社会关系与思想观念的论述、党和国家关于培育社会主义核心价值观和推进高校思想政治工作的政策文件，以及国内外社会学、传播学与思政教育学在弱关系、青年轻社交和价值观演变领域的丰硕成果，为本研究搭建分析框架与确立评判标准提供了充分支撑。

第二，网络语料来源真实充足。小红书、豆瓣高校小组、知乎高等教育板块及高校公开论坛汇集了大量在校大学生围绕搭子社交、求职备考互助、亲属往来和校园公共事务的真实记录，通过身份核验、情境限定与分层抽样，能够稳定建成3000至5000条有效文本数据库，满足宏观统计与人工深度编码的样本要求。

第三，研究方法与工具条件具备。本研究采用的文献分析、Python网络文本清洗与主题聚类以及NVivo人工内容编码均为社会科学领域成熟可靠的研究方法，前期已具备相关文献积累与文本数据处理基础，能够在既有准备之上顺利推进实证分析与对策研究。

\subsection{研究的创新点与重难点}

\subsubsection{研究可能存在的创新点}

第一，在研究视角上实现跨学科融通。既有研究多各自展开：社会学与传播学偏重考察弱连接在获取信息、促进求职或缓解焦虑上的实用功能，较少追问其对价值判断尺度的改变；思想政治教育学多侧重主流价值观的宏观阐释和网络不良思潮的内容剖析，较少剖析日常交往方式本身的影响。本研究回到马克思主义关于社会关系影响思想观念的根本立场，把弱连接这一日常交往形式作为切入点，系统剖析其微观作用过程如何冲击马克思主义价值观与社会主义核心价值观的培育，拓宽了思想政治教育学考察青年日常社会关系的研究视角。

第二，在研究内容上细化边界并打通关联场景。本研究以“对象可辨认、有实质回应、有一定延续性”收紧弱连接交往的界定边界，避免将其与匿名网络舆论混为一谈；依据公民道德生活领域确立边界清晰的四维价值观指标，从片面比较、低义务规则延伸、浅层满足替代与松散群体责任淡化四条微观路径揭示作用过程。尤其在人际交往维度，打破了将校园轻社交与家庭关系割裂考察的局限，将校园搭子社交与大学生疏离远亲的现象置于同一交往规则延伸的逻辑中加以阐释，并提出尊重个人边界与重构实体集体相结合的应对思路。

第三，在研究方法上加强微观经验支撑。针对既有思政教育研究对日常流动社交话语挖掘相对不足的局限，建立高校身份核验与弱连接交往情境双重限定的多平台网络文本采集流程，将场景词与价值词相分离的宏观统计分析同经过一致性检验的NVivo人工深度编码结合起来，使微观作用过程剖析获得贴近青年真实生活的经验依据。

\subsubsection{研究重点}

一是界定弱连接交往边界并搭建四维价值观分析框架。厘清弱连接交往与匿名网络舆论的区别，依据马克思主义价值立场、社会主义核心价值观与公民道德生活领域划分，确立考察大学生个人生活、人际交往、学习劳动与择业、社会公共四类价值观变化的清晰标准。

二是揭示负面影响表现、微观作用过程与深层成因。依托多平台大学生网络文本材料，系统呈现弱连接交往引发的四类价值观负面偏移，重点拆解四条微观传导路径并剖析背后的经济社会、家庭与高校育人成因。

三是构建契合高校思想政治教育规律的系统应对策略。围绕高校育人主责，结合青年重视个人边界的交往特点，从实体共同体建设、弱连接资源反向利用、思政课教学转化和校社协同四个方面提出可落地的引导方案。

\subsubsection{研究难点}

一是抽象价值规范向具体文本判断标准的转化。马克思主义价值观和社会主义核心价值观内涵宏大，而大学生在网络社区中的日常发言零散含蓄，还常夹杂反讽与调侃，需要制定边界清晰、操作性强且便于稳定判别的四维编码细则，确保人工编码的准确性与一致性。

二是开放网络社区中弱连接交往情境的精准筛选与解释限度把握。网络平台用户身份混杂，充斥着大量陌生人单向浏览、情绪宣泄和匿名争吵内容，需要在采集阶段同步核验高校在校生身份并严格限定双向互动情境，同时在阐释中严守网络文本资料的解释边界，避免作出超出经验材料支撑的推论。